% Claude 1962 proper guide, Twelfth Sunday after Pentecost: leaf-local
% typesetting primitives, shared by the canonical full edition (main.tex) and
% by the mechanical synthesis companion (synthesis.tex), which inputs it and
% then renews the edition notes. This file owns presentation only. Every
% textual, historical, exegetical, and rights judgment lives in the section
% files and the leaf's research records.
%
% Web-edition constraint: this leaf declares itself web-eligible, so it uses
% no landscape environment, no multicols, and no TikZ. Any ordering that a
% diagram might have carried is expressed as an ordered list or a table.

\usepackage{xurl}

\setlength{\emergencystretch}{2.4em}
\widowpenalty=10000
\clubpenalty=10000
\displaywidowpenalty=10000

\newcommand{\work}[1]{\textit{#1}}
\newcommand{\term}[1]{\textit{#1}}

% The missal accents the ligature in quǽsumus. U+01FD is outside the T1
% encoding's own repertoire, so it is set as an accented ligature rather than
% silently normalised to an unaccented "quæsumus".
\DeclareUnicodeCharacter{01FD}{\'{\ae}}
\DeclareUnicodeCharacter{01FC}{\'{\AE}}

% The missal's versicle and response signs are not in the T1 encoding. They
% are reproduced as a slashed V and R so that the appointed text can be
% transcribed as printed rather than silently normalised to "V." and "R.".
\usepackage{graphicx}
\newcommand{\Vsign}{\leavevmode\hbox{\rlap{\kern0.16em\raisebox{0.12ex}{\scalebox{0.85}[1.0]{/}}}V}}
\newcommand{\Rsign}{\leavevmode\hbox{\rlap{\kern0.16em\raisebox{0.12ex}{\scalebox{0.85}[1.0]{/}}}R}}
\DeclareUnicodeCharacter{2123}{\Vsign}
\DeclareUnicodeCharacter{211F}{\Rsign}
% The Gospel heading's Maltese cross, likewise outside T1.
\newcommand{\gospelcross}{\ensuremath{\boldsymbol{+}}}
\usepackage{amsmath,amsbsy}
\DeclareUnicodeCharacter{2720}{\gospelcross}

% Quiet parenthetical proper cues, in liturgical order, for later headings and
% cross-references. Defined once so the abbreviations stay uniform.
\newcommand{\cue}[1]{\textnormal{(\textit{#1})}}

% Keeps an element's heading, identification line and English block together.
\newcommand{\elementguard}{\Needspace{12\baselineskip}}

% Element identification line: Latin incipit, printed heading and reference,
% and the missal's marginal number.
\newcommand{\elementid}[3]{%
  \par{\small\latin{#1} \textperiodcentered{} #2 \textperiodcentered{} marginal no.~#3}\par}

% The canonical full edition prints the complete appointed Latin.
\newenvironment{appointedlatin}[1]{%
  \begin{massbox}{#1}\itshape}{\end{massbox}}
\newcommand{\appointedrubric}[2]{\par\textbf{Rubric as printed:} \latin{#1}
\textperiodcentered{} #2\par}

% An identified public-domain English witness, never a project translation.
% The title always names the witness and the exact locus it was taken from.
\newenvironment{namedtranslation}[1]{%
  \begin{sourcecard}{#1}\small}{\end{sourcecard}}

% A place where the registered public-domain English does not answer the
% appointed Latin. The gap is stated, never filled.
\newcommand{\englishgap}[1]{\par\textbf{Where the registered English does not
answer the Latin.} #1\par}

% The two editions share every section file below. These two passages are the
% only reader-facing text that differs between them, so each edition's
% entrypoint owns its own wording and the sections stay identical. The full
% edition's wording is defined here; synthesis.tex renews both.
\newcommand{\latinpolicynote}{The complete appointed Latin has already been
printed above, in \textit{The Appointed Formulary in Full}, in liturgical order
and beside the same English; the commentary quotes it again only where an
argument turns on the wording.}
\newcommand{\editionnote}{\textbf{Which edition this is.} This canonical
\textit{full edition} carries the map, the dossiers and the thematic movement,
then the complete appointed formulary --- the Latin of the 1962 Vatican
\latin{editio typica} with the English of the registered public-domain
witnesses --- and then the element-by-element commentary, synthesis, and
apparatus from the same source leaf. The mechanical \textit{synthesis edition}, published at the same
identity with a \texttt{-synthesis} suffix, omits the complete appointed texts
and the element-by-element sweep and carries instead one integrated
commentary; it adds no research of its own. The two editions rest on one body
of research and share one set of source records.}

% A witness dossier inside the detailed commentary: who, at what exact locus,
% arguing what. The left rule keeps it distinct from surrounding synthesis
% without implying doctrinal weight.
\newenvironment{witnessnote}{\begin{studyframe}\small}{\end{studyframe}}
\newcommand{\wlocus}[1]{\textbf{#1}\quad}

% Governing thesis for the thematic movement. Substantive content, permitted
% to open the body of its own section.
\newcommand{\movementthesis}[1]{%
  \begin{studybox}{Governing thesis}#1\end{studybox}}

% The exploratory-proposal notice and the proposals themselves. Every proposal
% renders its anchors, mechanism, fruit, and controlling limit as labelled
% fields so no reader can mistake a proposal for a sourced finding.
\newenvironment{proposal}[1]{%
  \begin{massbox}{Editorial proposal --- #1}\small}{\end{massbox}}
\newcommand{\pfield}[2]{\par\textbf{#1:} #2}

% Notable-and-quotable gallery entry: appointed phrase, later user, locus,
% and the turn in register. Breakable, so a long entry splits rather than
% jumping whole to the next page and leaving a sparse one behind.
\newenvironment{afterlife}[1]{%
  \begin{massbox}{#1}\small}{\end{massbox}}

% Two-tier page-2 dossier table, at the measurements the profile fixes.
\newenvironment{dossiertable}
{\footnotesize
 \renewcommand{\arraystretch}{1.02}%
 \setlength{\LTpre}{0.15em}\setlength{\LTpost}{0pt}%
 \begin{longtable}{@{}>{\raggedright\arraybackslash}p{0.14\linewidth}>{\raggedright\arraybackslash}p{0.18\linewidth}>{\raggedright\arraybackslash}p{0.35\linewidth}>{\raggedright\arraybackslash}p{0.16\linewidth}@{}}
 \toprule
 \textbf{Proper} & \textbf{Citation} & \textbf{Location} & \textbf{Date}\\
 \midrule
 \endhead}
{\bottomrule\end{longtable}}

% Full-width explanatory row inside a dossier.
\newcommand{\dossierprose}[1]{\multicolumn{4}{@{}>{\raggedright\arraybackslash}p{0.92\linewidth}@{}}{#1}\\[0.1em]}
\newcommand{\dossierevent}[1]{\multicolumn{4}{@{}>{\raggedright\arraybackslash}p{0.92\linewidth}@{}}{\textit{#1}}\\}

% Reception matrix and comparison tables used in the detailed commentary.
\newenvironment{receptiontable}[3]
{\footnotesize\begin{longtable}{>{\raggedright\arraybackslash}p{0.20\linewidth}>{\raggedright\arraybackslash}p{0.34\linewidth}>{\raggedright\arraybackslash}p{0.38\linewidth}}
\toprule
\textbf{#1} & \textbf{#2} & \textbf{#3}\\
\midrule
\endhead}
{\bottomrule\end{longtable}}

% Attributions the guide corrects: the claim a reasonable guide could make,
% what the checked sources actually show, and the proper the evidence sits at.
\newenvironment{correctionstable}
{\footnotesize\begin{longtable}{>{\raggedright\arraybackslash}p{0.26\linewidth}>{\raggedright\arraybackslash}p{0.52\linewidth}>{\raggedright\arraybackslash}p{0.12\linewidth}}
\toprule
\textbf{A claim a guide could make} & \textbf{What the checked sources show} & \textbf{Proper}\\
\midrule
\endhead}
{\bottomrule\end{longtable}}

% \sectionguard and the deepstudy reading size come from the shared preamble
% and are deliberately not redefined here.
